\documentclass[11pt,a4paper]{article}

% ---------- Packages ----------
\usepackage[a4paper,margin=1.5cm]{geometry}
\usepackage[T1]{fontenc}
\usepackage[utf8]{inputenc}
\usepackage{lmodern}
\usepackage{microtype}
\usepackage{xcolor}
\usepackage[hidelinks]{hyperref}
\usepackage{tabularx}
\usepackage{enumitem}
\usepackage{titlesec}
\usepackage{ragged2e}
\usepackage{fontawesome5}

% ---------- Styling ----------
\definecolor{Accent}{HTML}{1F4E79}
\definecolor{Muted}{HTML}{6B7280}

\setlength{\parindent}{0pt}
\setlength{\parskip}{0.3em}

\titleformat{\section}
  {\large\bfseries\color{Accent}}
  {}{0pt}{}
  [\vspace{-0.1em}\titlerule]

\titlespacing*{\section}{0pt}{0.9em}{0.5em}

\setlist[itemize]{leftmargin=1.2em, topsep=0.25em, itemsep=0.15em, parsep=0em}

% ---------- Personal Information ----------

\newcommand{\FullName}{Arseniy Obolenskiy}
\newcommand{\Email}{good.doog.2012@gmail.com}
\newcommand{\Location}{Munich, Bavaria, Germany}

\newcommand{\LinkedInURL}{https://www.linkedin.com/in/arseniy-obolenskiy/}
\newcommand{\GitHubURL}{https://github.com/aobolensk}

% ---------- Helpers ----------
\newcommand{\iconsep}{\hspace{0.6em}}
\newcommand{\muted}[1]{{\color{Muted}#1}}

\newcommand{\cvheader}{%
  \begin{flushright}
    {\Huge\bfseries \FullName}\par
    \vspace{0.2em}
    \small
    \faEnvelope\iconsep
    \href{mailto:\Email}{\Email}%
    \iconsep
    \faMapMarker*\iconsep\Location\par
    \vspace{0.15em}
    \small
    \faLinkedin\iconsep\href{\LinkedInURL}{LinkedIn}\iconsep
    \faGithub\iconsep\href{\GitHubURL}{GitHub}%
    \par
    \vspace{0.6em}
  \end{flushright}
}

\newcommand{\cventry}[5]{%
  \par\noindent
  \begin{tabularx}{\linewidth}{@{}p{0.17\linewidth}@{\hspace{0.55em}}X@{}}
    \muted{\small #1} &
    {\textbf{#3} \hfill \muted{\small #4}}\\[-0.15em]
    & {\textit{#2}}\\
  \end{tabularx}\par
  #5
  \vspace{0.25em}
}

\newcommand{\cvbullets}[1]{%
  \begin{itemize}[leftmargin=1.4em, labelsep=0.5em]
    #1
  \end{itemize}
}

\newcommand{\cvplain}[1]{%
  \vspace{0.1em}
  \small #1
  \par
}

% ---------- Document ----------
\begin{document}
\cvheader

% ---------- Summary ----------
\section*{Profile}
\justifying
I am a software engineer working in compiler development industry. I have experience in C++ and Python programming, worked in projects related to media software stack and currently focused on working in compiler development and AI frameworks scope. Furthermore, I am participating in knowledge sharing activities, involved into teaching and mentoring. Also, I am fond of doing pet open source projects in my spare time on GitHub to learn some technologies besides professional job experience.
I am looking for professional growth and strive to learn something new and develop skills.

% ---------- Work Experience ----------
\section*{Work Experience}

\cventry
  {03.2026 -- present}
  {Advanced Micro Devices (AMD)}
  {Senior Member of Technical Staff}
  {Aschheim, Bavaria, Germany}
  {\cvbullets{
    \item Working on LLVM compiler (mostly SPIR-V backend, MLIR SPIR-V dialect and AMDGPU backend) and ROCm stack
    \item Related open source repository: \href{https://github.com/llvm/llvm-project}{github.com/llvm/llvm-project}
  }}

\cventry
  {01.2025 -- 02.2026}
  {Intel Corporation}
  {AI Frameworks Engineer}
  {Feldkirchen, Bavaria, Germany}
  {\cvbullets{
    \item Enabled acceleration of common LLM subgraph patterns on model compilation step, gaining performance for multiple model variants on both x86, ARM and RISC-V platforms
    \item Maintained RISC-V CPU inference component, drove collaboration activities for new outside contributors
    \item Led code-quality improvement initiatives by adopting clang-tidy static analysis across a large codebase, performing extensive code cleanup, and enforcing best practices reducing overall code review time
    \item Related open source repository: \href{https://github.com/openvinotoolkit/openvino}{github.com/openvinotoolkit/openvino}
  }}

\cventry
  {09.2025 -- 06.2026}
  {Nizhny Novgorod State University}
  {Research Supervisor}
  {Remote (contract)}
  {%
    \cvplain{Supervised the research and thesis projects of three Bachelor's students. As a result, all three theses received the highest possible grade}
  }

\cventry
  {02.2024 -- 06.2026}
  {Nizhny Novgorod State University}
  {Lecturer}
  {Remote (contract)}
  {%
   \cvplain{Worked as a lecturer for several subjects: \textit{Introduction to Compilers}, \textit{Parallel Programming}, \textit{Basics of AI accelerators programming}. Conducted practical classes, reviewed student assignments, and automated homework verification.}
   \vspace{0.1em}
   \textbf{Introduction to Compilers}\par
   \cvbullets{
     \item Extended and updated materials for ``Introduction to Compilers'' course
     \item Delivered lectures on ``Introduction to Compilers''
     \item Introduced short tests component for the lectures to measure subject understanding among the students
     \item Prepared an infrastructure for practical component of the course and participated in student laboratory assignments review
     \item Prepared publicly available course materials: \href{https://github.com/NN-complr-tech/Complr-course-lectures}{github.com/NN-complr-tech/Complr-course-lectures}
   }
   \par
   \textbf{Parallel Programming}\par
   \cvbullets{
     \item Delivered practice oriented lectures as a part of ``Parallel Programming'' course
     \item Prepared publicly available course materials: \href{https://github.com/learning-process/parallel_programming_slides}{github.com/learning-process/parallel\_programming\_slides}
   }
   \par
   \textbf{Basics of AI accelerators programming}\par
   \cvbullets{
     \item Implemented part of lectures about Ascend NPU from scratch
     \item Delivered lectures on ``Basics of AI accelerators programming'' about Ascend NPU
   }
  }

\cventry
  {09.2023 -- 12.2024}
  {Huawei}
  {Senior Software Engineer}
  {Dongguan, Guangdong, China / Suzhou, Jiangsu, China}
  {\cvbullets{
    \item Worked on MLIR based solution for HW specific code optimizations for ML operators written for different ML frameworks, implementing and enabling source-to-source compiler based transformation
    \item Designed and implemented an end-to-end translator for OpenAI Triton programs targeting Ascend NPU devices
    \item Worked on binary profiling tool implementation using binary instrumentation technology on LLVM stack and custom written from scratch non-LLVM infrastructure
    \item Prototyped conversion tool for source-to-source transformation of OpenMP pragma directives using clang frontend capabilities
    \item Investigated different variation of compiler based source-to-source code transformations for conversion between different programming models
    \item Performed team research project deliverables' integration into headquarters product mainline, meeting and negotiating with headquarters team members; successfully integrated two projects developed by overseas team into the product mainline
  }}

\cventry
  {09.2022 -- 09.2023}
  {Huawei}
  {Software Engineer}
  {Nizhny Novgorod, Russia}
  {\cvbullets{
    \item Worked on LLVM-based compiler for NPU
    \item Implemented whole compilation stack from FE programming model to backend code generation: clang, MLIR, LLVM, target specific assembly
    \item Implemented solution for source-to-source MLIR based code transformation for code conversion between different programming models (SIMT and SIMD)
    \item Implemented binary profiling tool using binary instrumentation technology on LLVM stack
    \item Was in charge of CI infrastructure for all projects in research team
  }}

\cventry
  {12.2021 -- 09.2022}
  {Intel Corporation}
  {Compiler Engineer}
  {Nizhny Novgorod, Russia}
  {\cvbullets{
    \item Working on various features and fixing bugs reported by the community in DPC++ compiler
    \item Supporting SYCL device headers
  }}

\cventry
  {08.2021 -- 12.2021}
  {Intel Corporation}
  {Graphics Software Engineer}
  {Nizhny Novgorod, Russia}
  {\cvbullets{
    \item Worked on a PoC of a new SDK for media stack related to screen sharing and broadcasting
  }}

\cventry
  {12.2020 -- 12.2022}
  {Nizhny Novgorod State University}
  {Laboratory Research Assistant (part-time)}
  {Nizhny Novgorod, Russia}
  {\cvbullets{
    \item Worked on graph and ML algorithms application for solving biological problems
    \item Worked on researching OpenVINO neural network layers optimization for ARM CPU architecture
  }}

\cventry
  {01.2020 -- 08.2021}
  {Intel Corporation}
  {Graphics Software Intern}
  {Nizhny Novgorod, Russia}
  {\cvbullets{
    \item Worked on various features and supported MediaSDK and oneVPL products for Intel GPUs
    \item Worked with customers and community
    \item Supported C++ validation framework
    \item Related open source repositories:
      \begin{itemize}
        \item \href{https://github.com/Intel-Media-SDK/MediaSDK}{github.com/Intel-Media-SDK/MediaSDK}
        \item \href{https://github.com/oneapi-src/oneVPL-intel-gpu}{github.com/oneapi-src/oneVPL-intel-gpu}
      \end{itemize}
  }}

% ---------- Volunteering ----------
\section*{Volunteering}

\cventry
  {10.2021 -- 06.2026}
  {Nizhny Novgorod State University}
  {Mentor and Project Supervisor (ITLab)}
  {Remote}
  {\cvbullets{
    \item Participated in ITLab project organized by university to drive additional voluntary IT classes for Bachelor's degree students. These classes are dedicated to a certain topic where the project is being implemented by the team of 2-6 students.
    \item Topics:
      \begin{itemize}
        \item Oct 2025 - Jun 2026 (6 students): Android applications development with LLM based on the OpenVINO Java API and ARM CPU plugin (students earned second place in the university project competition)
        \item Oct 2025 - Jun 2026 (3 students): Optimization of neural network layers for LLM tasks on embedded architectures (RISC-V) in OpenVINO Toolkit (students earned second place in the university project competition)
        \item Oct 2024 - Jun 2025 (4 students): Neural network inference using ARM Compute Library and industrial software development
        \item Oct 2023 - Jun 2026 (3-4 students): Neural network inference framework implementation and industrial software development (3 students successfully developed this project into their bachelor's theses)
        \item Oct 2022 - Jun 2023 (2 students): Computer vision algorithms optimization on embedded devices
        \item Oct 2021 - Jun 2022 (2 students): Neural networks optimization on ARM architecture
      \end{itemize}
    \item Most of the projects and related materials are available on GitHub: \href{https://github.com/embedded-dev-research}{https://github.com/embedded-dev-research}
  }}

% ---------- Education ----------
\section*{Education}

\cventry
  {10.2023 -- 07.2026}
  {Nizhny Novgorod State University}
  {PhD Program, Computer Science and Informatics}
  {Nizhny Novgorod, Russia}
  {\cvplain{PhD candidate: Coursework completed, ABD (all but dissertation)}}

\cventry
  {09.2021 -- 06.2023}
  {Nizhny Novgorod State University}
  {Master's degree with honors}
  {Nizhny Novgorod, Russia}
  {\cvbullets{
    \item Field of study: Fundamental Informatics and Information Technologies
    \item Neural network optimization research on ARM architecture
    \item Was a mentor for a group of students in ITLab project
    \item Final qualification work: ``Neural networks inference optimization on ARM architecture''
  }}

\cventry
  {09.2017 -- 06.2021}
  {Nizhny Novgorod State University}
  {Bachelor's degree with honors, Software Engineering}
  {Nizhny Novgorod, Russia}
  {\cvbullets{
    \item Graph algorithms and data structures research; graph coarsening implementation; neural networks application to tasks related to graphs
    \item Participated in ITLab project
    \item Final qualification work: ``Building graphs that are topologically similar to the original ones to improve the accuracy of machine learning algorithms''
  }}

% ---------- Publications ----------
\section*{Publications}
\cventry
  {02.01.2026}
  {Springer, Supercomputing (RuSCDays 2025), LNCS vol. 16196}
  {Practical Aspects of Teaching Parallel Programming at the Lobachevsky University}
  {}
  {%
    \cvplain{\textbf{Authors:} Alexander Nesterov, \underline{Arseniy Obolenskiy}, Alexander Sysoyev, Iosif Meyerov}
    \cvplain{\textbf{Language:} English}
    \cvplain{\textbf{DOI:} \href{https://doi.org/10.1007/978-3-032-13127-0_37}{\nolinkurl{10.1007/978-3-032-13127-0_37}} \quad
    }
  }

\cventry
  {19.11.2025}
  {Mathematical Modeling and Supercomputer Technologies Conference 2025}
  {Collaborative development of a library for convolutional neural networks inference}
  {}
  {%
    \cvplain{\textbf{Authors:} A. A. Sorokin, S. M. Titov, A. Yu. Nesterov, \underline{A. A. Obolenskiy}}
    \cvplain{\textbf{Language:} Russian}
    \cvplain{\textbf{Link:} \href{https://web.archive.org/web/20251209131202id_/https://mmst.unn.ru/wp-content/uploads/2025/12/MMST2025_Proceedings.pdf\#page=152}{\nolinkurl{mmst.unn.ru/wp-content/uploads/2025/12/MMST2025_Proceedings.pdf\#page=152}} \quad
    }
  }

\cventry
  {14.11.2022}
  {Mathematical Modeling and Supercomputer Technologies Conference 2022}
  {Neural network layers optimization on ARM architecture}
  {}
  {%
    \cvplain{\textbf{Authors:} \underline{A. A. Obolenskiy}, A. Yu. Nesterov, I. B. Meyerov}
    \cvplain{\textbf{Language:} Russian}
    \cvplain{\textbf{Link:} \href{https://web.archive.org/web/20221124173857id_/https://hpc-education.unn.ru/files/conference_hpc/2022/MMST2022_Proceedings.pdf\#page=86}{\nolinkurl{hpc-education.unn.ru/files/conference_hpc/2022/MMST2022_Proceedings.pdf\#page=86}} \quad
    }
  }

\cventry
  {22.11.2021}
  {Mathematical Modeling and Supercomputer Technologies Conference 2021}
  {Application of machine learning methods to determine the origin of people based on the results of a mutation in their DNA}
  {}
  {%
    \cvplain{\textbf{Authors:} \underline{A. A. Obolenskiy}, V. O. Devlikamov, A. I. Kalyakulina, A. Yu. Nesterov, I. B. Meyerov, M. V. Ivanchenko}
    \cvplain{\textbf{Language:} Russian}
    \cvplain{\textbf{Link:} \href{https://web.archive.org/web/20220331185330id_/https://hpc-education.unn.ru/files/conference_hpc/2021/MMST2021_Proceedings.pdf\#page=251}{\nolinkurl{hpc-education.unn.ru/files/conference_hpc/2021/MMST2021_Proceedings.pdf\#page=251}} \quad
    }
  }

\cventry
  {21.09.2020}
  {International Conference Russian Supercomputing Days 2020}
  {Comparison of graph coarsening algorithms in parallel sparse matrix reordering task}
  {}
  {%
    \cvplain{\textbf{Authors:} A. Yu. Pirova, \underline{A. A. Obolenskiy}, V. V. Devlikamov, A. Yu. Nesterov}
    \cvplain{\textbf{Language:} Russian}
    \cvplain{\textbf{Link:} \href{https://web.archive.org/web/20201110231956id_/https://2020.russianscdays.org/files/2020/RuSCDays20_Proceedings.pdf\#page=142}{\nolinkurl{2020.russianscdays.org/files/2020/RuSCDays20_Proceedings.pdf\#page=142}} \quad
    }
  }

% ---------- Languages & Skills ----------
\section*{Languages}
\begin{tabular}{@{}l@{\hspace{0.8em}}l@{}}
  \textbf{English} & upper intermediate (B2) \\
  \textbf{Russian} & native \\
  \textbf{German} & intermediate (B1, Goethe-Zertifikat) \\
  \textbf{Mandarin Chinese} & elementary (HSK 2) \\
\end{tabular}

\section*{Professional Skills}
\small
C/C++, %
Python, %
LLVM, %
MLIR, %
source-to-source transformations, %
code generation, %
performance optimization and profiling, %
static code analysis, %
CI/CD, %
AI frameworks (PyTorch, OpenVINO), %
LLM inference optimization, %
parallel programming (OpenMP, SYCL), %
CPU/NPU architecture-aware optimization (x86, ARM, RISC-V), %
binary instrumentation, %
Linux, %
Git, %
research project leadership, %
university lecturing and mentoring

\end{document}
